\section*{Chapitre \ref{ch:b1:newton-dynamics} --- Les lois de Newton}

\begin{solution}{exo:b1:newton-dynamics:1}
Indication du pèse-personne $N = m(g + a)$ avec $a$ l'accélération vers
le haut~: vers le haut, $70 \times 11.81 = \qty{827}{N}$ («~\qty{84}{kg}~»)~; vers le bas, $70 \times 7.81 =
\qty{547}{N}$ («~\qty{56}{kg}~»)~; à vitesse constante, \qty{687}{N}
(\qty{70}{kg})~; en chute libre, $0$.
\end{solution}

\begin{solution}{exo:b1:newton-dynamics:2}
Même accélération $a$ pour les deux (fil inextensible)~: $m_1a = T$,
$m_2a = m_2g - T$, donc $a = m_2g/(m_1 + m_2) = \qty{3.3}{m/s^2}$ et
$T = m_1a = \qty{6.5}{N}$.
\end{solution}

\begin{solution}{exo:b1:newton-dynamics:3}
$a = g(\sin\alpha - \mu_d\cos\alpha) = 9.81(0.50 - 0.17) = \qty{3.2}{m/s^2}$~;
$v = \sqrt{2 \times 3.2 \times 5.0} = \qty{5.7}{m/s}$. Aucune \omterm{def:b1:newton-dynamics:momentum}{masse} en vue.
\end{solution}

\begin{solution}{exo:b1:newton-dynamics:4}
$k\Delta\ell = mg$~: $\Delta\ell = 0.50 \times 9.81/200 = \qty{2.5}{cm}$.
Avec $x$ l'écart à l'équilibre, $m\ddot x = -kx$ (le poids et
l'allongement statique se compensent)~: $\omega_0 = \sqrt{k/m} = \qty{20}{rad/s}$,
$T = \qty{0.31}{s}$.
\end{solution}

\begin{solution}{exo:b1:newton-dynamics:5}
$\mu_s = \tan 22^\circ = 0.40$. Virage~: $v \leq \sqrt{\mu_sgR}$~: $\sqrt{0.8
\times 9.81 \times 50} = \qty{20}{m/s}$ (\qty{71}{km/h})~; sur route
mouillée, \qty{14}{m/s} (\qty{50}{km/h}).
\end{solution}

\begin{solution}{exo:b1:newton-dynamics:6}
Parachutiste~: $v_\infty = \sqrt{2 \times 785/(1.2 \times 0.70)} = \qty{43}{m/s}$~;
$\mathrm{Re} = \rho vL/\eta \approx 1.2 \times 43 \times 0.5/1.8\times10^{-5}
\approx 10^6$~: quadratique. Gouttelette~: $m = \tfrac43\pi r^3\rho_w = \qty{4.2e-12}{kg}$,
$v_\infty = mg/6\pi\eta r = 4.1\times10^{-11}/3.4\times10^{-9} = \qty{1.2}{cm/s}$~;
$\mathrm{Re} = 1.2 \times 0.012 \times 2\times10^{-5}/1.8\times10^{-5} = 0.02$~:
linéaire.
\end{solution}

\begin{solution}{exo:b1:newton-dynamics:7}
En bas~: $v^2 = 2g\ell(1 - \cos 60^\circ) = g\ell$~; $T = mg + mv^2/\ell =
2mg$. Au point de lâcher~: $v = 0$, $T = mg\cos\theta_0 = 0.5mg$.
\end{solution}

\begin{solution}{exo:b1:newton-dynamics:8}
$v = \sqrt{gR\tan\beta} = \sqrt{9.81 \times 200 \times 0.268} = \qty{23}{m/s}$
(\qty{82}{km/h}). Plus lentement~: la voiture tend à glisser vers le bas
du dévers, le frottement doit pointer vers le haut de la pente~; plus
vite~: il pointe vers le bas pour ajouter de la \omterm{prop:b1:newton-dynamics:centripetal}{force centripète}.
\end{solution}

\begin{solution}{exo:b1:newton-dynamics:9}
$T\cos\alpha = mg$, $T\sin\alpha = m\omega^2\ell\sin\alpha$~: $\omega =
\sqrt{g/\ell\cos\alpha} = \sqrt{9.81/0.866} = \qty{3.4}{rad/s}$, période
\qty{1.9}{s}~; $T = mg/\cos\alpha = 1.15mg$.
\end{solution}

\begin{solution}{exo:b1:newton-dynamics:10}
$a = 27.8/8.0 = \qty{3.5}{m/s^2}$, $F = \qty{3.5}{kN}$~; $F \leq \mu_smg$
donne $\mu_s \geq a/g = 0.35$. Sur route mouillée, $\mu_s \approx 0.4$
plafonne l'accélération vers $0.4g$ quel que soit le moteur~:
c'est l'adhérence, non la puissance.
\end{solution}

\begin{solution}{exo:b1:newton-dynamics:11}
$v = v_\infty(1 - \eu^{-t/\tau})$, $v_\infty = g\tau = \qty{4.9}{m/s}$~;
$z = \int v = v_\infty[t - \tau(1 - \eu^{-t/\tau})]$. À \qty{2.0}{s}~:
$4.9 \times (2.0 - 0.5 \times 0.98) = \qty{7.4}{m}$, contre $\tfrac12 gt^2
= \qty{20}{m}$ dans le vide.
\end{solution}

\begin{solution}{exo:b1:newton-dynamics:12}
Projection radiale (vers le centre)~: $N - mg\cos\theta = mv^2/R$, donc
$N = m(v^2/R + g\cos\theta)$. Au sommet ($\theta = \pi$)~: $N = m(v^2/R - g)
\geq 0$ si et seulement si $v \geq \sqrt{gR}$. Avec la loi des vitesses, $v_{\text{sommet}}^2 =
v_0^2 - 4gR \geq gR$~: $v_0 \geq \sqrt{5gR}$~; alors, en bas, $N =
m(5g + g) = 6mg$.
\end{solution}

\begin{solution}{pb:b1:newton-dynamics:1}
\textbf{1.} $m\dot v = mg$~: $v = gt$, $z = \tfrac12 gt^2$~; à \qty{10}{s}~:
\qty{98}{m/s} et \qty{490}{m}.

\textbf{2.} $mg = \qty{785}{N}$.

\textbf{3.} Sur le sol, le plancher pousse vers le haut avec $N = mg$ ---
c'est cette poussée que l'on ressent comme son poids. En chute libre, il
n'y a aucune \omterm{def:b1:newton-dynamics:momentum}{force} de contact~: la pesanteur agit également sur chaque
partie du corps et rien n'appuie.

\textbf{4.} Oui, à quelques dixièmes de pour cent près sur une minute~;
sont négligées~: la rotation de la Terre (termes d'inertie d'entraînement
et de Coriolis, \cref{ch:b1:non-inertial-frames}) et le vent.

\textbf{5.} $F_d = \tfrac12\rho C_xSv^2 = 0.48\,v^2$ (unités SI).

\textbf{6.} $m\,\dd v/\dd t = mg - 0.48v^2$.

\textbf{7.} $v_\infty = \sqrt{mg/0.48} = \sqrt{1635} = \qty{40}{m/s} =
\qty{146}{km/h}$.

\textbf{8.} En divisant par $mg$~: $\dd u/\dd t\,(v_\infty/g) = 1 - u^2$~;
$\tau = 40.4/9.81 = \qty{4.1}{s}$.

\textbf{9.} $\int_0^u\dd u'/(1 - u'^2) = \operatorname{artanh}u = t/\tau$,
donc $u = \tanh(t/\tau)$.

\textbf{10.} $t = \tau\operatorname{artanh}(0.9) = 1.47\tau = \qty{6.1}{s}$~;
pour $0.99$~: $2.65\tau = \qty{11}{s}$.

\textbf{11.} $z = \int v_\infty\tanh(t/\tau)\dd t = v_\infty\tau\ln\cosh(t/\tau)$~;
à \qty{10}{s}~: $166 \times \ln\cosh(2.43) = 166 \times 1.74 = \qty{290}{m}$,
contre \qty{490}{m}.

\textbf{12.} $\mathrm{Re} = 1.2 \times 40 \times 0.5/1.8\times10^{-5} \approx
\num{1.3e6}$~: bien au-dessus de $1000$, donc quadratique.

\textbf{13.} $v_\infty' = \sqrt{785/(0.6 \times 25)} = \sqrt{52} = \qty{7.2}{m/s}$.

\textbf{14.} Traînée $= 15 \times 1600 = \qty{24}{kN}$~; $a = (24000 - 785)/80
= \qty{290}{m/s^2} \approx 30g$~: le harnais casserait, ou le sauteur.

\textbf{15.} En moyenne $a \approx (40 - 7)/3 = \qty{11}{m/s^2} \approx 1.1g$~;
\omterm{def:b1:newton-dynamics:momentum}{force} du harnais $\approx m(g + a) \approx \qty{1.7}{kN}$.

\textbf{16.} $\tau' = 7.2/9.81 = \qty{0.73}{s}$~: quelques secondes après
l'ouverture complète, la descente est établie.

\textbf{17.} $\vect a = \vect 0$~: \omterm{def:b1:newton-dynamics:momentum}{force} du harnais $= mg = \qty{785}{N}$.

\textbf{18.} L'équilibre $mg = \tfrac12\rho C_xSv^2$ donne $v \propto (C_xS)^{-1/2}$~:
avec la moitié de l'aire, la vitesse est multipliée par $\sqrt2$ (\qty{10}{m/s}).

\textbf{19.} $v \propto \sqrt m$~: $7.2 \times \sqrt{100/80} = \qty{8.1}{m/s}$.

\textbf{20.} $h = v^2/2g = 52/19.6 = \qty{2.6}{m}$~: un saut depuis un
premier étage.

\textbf{21.} $a = v^2/2d$~: \qty{52}{m/s^2} ($5.3g$) sur \qty{0.5}{m},
\omterm{def:b1:newton-dynamics:momentum}{force} $m(g + a) \approx \qty{5}{kN}$~; \qty{17}{m/s^2} ($1.8g$) sur
\qty{1.5}{m}, \qty{2.2}{kN}. Il faut rouler.

\textbf{22.} Une traînée plus grande (et un peu de portance) dépasse
momentanément le poids~: le sauteur ralentit et touche le sol
au-dessous de $v_\infty'$.

\textbf{23.} $m = \tfrac43\pi r^3\rho_w = \qty{4.2e-6}{kg}$, $S = \pi r^2 =
\qty{3.1e-6}{m^2}$~: $v_\infty = \sqrt{2mg/\rho C_xS} = \sqrt{8.2\times10^{-5}/
1.9\times10^{-6}} = \qty{6.6}{m/s}$~; $\mathrm{Re} = 1.2 \times 6.6 \times
2\times10^{-3}/1.8\times10^{-5} \approx 900$~: quadratique, à peu près.

\textbf{24.} $v_\infty = \qty{1.2}{cm/s}$ (\cref{exo:b1:newton-dynamics:6})~;
$\mathrm{Re} \approx 0.02 \ll 1$~: la \omterm{prop:b1:newton-dynamics:drag}{loi de Stokes} s'applique, de façon cohérente.

\textbf{25.} $m\dot{\vect v} = m\vect g + \vect F_{\text{traînée}}$~; la
traînée est linéaire en $v$ pour les objets petits et lents, quadratique
pour les gros et rapides. Les nombres~: \qty{40}{m/s} pour un corps en
chute, \qty{7}{m/s} sous une voilure (ou pour une goutte de pluie),
\qty{1}{cm/s} pour le brouillard --- la même loi sur quatre \omterm{def:b1:units-dimensions:oom}{ordres de grandeur}.
\end{solution}
